\chapter{Group Cohomology}

\begin{center}\itshape
When a group acts on an abelian group, three small cohomology groups answer three old questions: which elements are fixed, which ``twisted homomorphisms'' are trivial, and which group extensions exist. The answers
include Hilbert's Theorem~90, the Schur--Zassenhaus theorem, and the carry digit of school arithmetic.
\end{center}

Add $27$ and $45$ the way a child does. The units give $7+5=12$: write $2$, carry $1$. The tens give $2+4+1=7$. The carry is the one non-obvious step, and children struggle with it for good reason, as Daniel Isaksen
explained in 2002 in the \emph{American Mathematical Monthly}: the carry is a \emph{cocycle}. Write a number below $100$ as a pair of digits $(a,x)$, tens and units, each in $\Z/10$. Then
\[
(a,x)+(b,y)=\bigl(a+b+c(x,y),\ x+y\bigr),\qquad c(x,y)=\begin{cases}1&\text{if }x+y\ge10,\\0&\text{otherwise,}\end{cases}
\]
with all digit arithmetic modulo $10$. The function $c$ glues two copies of $\Z/10$ into $\Z/100$, and no clever choice of digit symbols can make the carry disappear: if it could, $\Z/100$ would be $\Z/10\times\Z/10$, in which every
element has order dividing $10$. This chapter shows that the carry is a non-trivial element of a group $H^2(\Z/10,\Z/10)\cong\Z/10$, which classifies all the ways of building a group of order $100$ from two copies of $\Z/10$ (one of them is $\Z/10\times\Z/10$, four are copies of $\Z/100$, and the others are $\Z/20\times\Z/5$ and $\Z/50\times\Z/2$).

The same group $H^2$ was met long before it had a name. Otto Hölder in 1893 and Otto Schreier in 1926 described all extensions of one group by another using ``factor systems'', which are cocycles. Issai Schur used them
in 1904 for projective representations. And in 1897, in his \emph{Zahlbericht}, the report on algebraic number theory that the German Mathematical Society had commissioned from him, David Hilbert proved as his ninetieth
theorem a statement about norms in cyclic extensions of number fields. ``Theorem~90'' kept its number, and became, in the hands of Emmy Noether and Andreas Speiser, the statement that a first cohomology group vanishes.
Cohomology of groups as a theory came in the 1940s, from Heinz Hopf's work in topology and from Eilenberg and Mac Lane; in Chapter~6's language $H^n(G,A)=\operatorname{Ext}^n_{\Z[G]}(\Z,A)$. This chapter needs only degrees $0$, $1$ and $2$, written by hand.

\section{Cochains and the first cohomology groups}

Let $G$ be a group. A \emph{$G$-module} is an abelian group $A$, written additively, with an action of $G$ by automorphisms: $g\cdot(a+b)=g\cdot a+g\cdot b$, $(gh)\cdot a=g\cdot(h\cdot a)$, $1\cdot a=a$. Equivalently, a module over the
group ring $\Z[G]$. Examples: any abelian group with the \emph{trivial} action; $\Z$ with $\Z/2$ acting by $\pm1$; for a Galois extension $L/K$ with group $G$, the additive group $L$ and the multiplicative group $L^\times$ (written
multiplicatively, of course); for $A\trianglelefteq E$ abelian with $E/A=G$, the group $A$ with $G$ acting by conjugation (Section~\ref{sec:extensions}).

\begin{definition}\label{def:cochains}
The \emph{$n$-cochains} are the functions $f\colon G^n\to A$ (for $n=0$, the elements of $A$). The \emph{coboundaries} $d\colon C^n\to C^{n+1}$ are, for $n=0,1,2$,
\[
\begin{gathered}
(da)(g)=g\cdot a-a,\qquad(df)(g,h)=g\cdot f(h)-f(gh)+f(g),\\
(df)(g,h,k)=g\cdot f(h,k)-f(gh,k)+f(g,hk)-f(g,h).
\end{gathered}
\]
A direct computation, which the reader should do once, gives $d\circ d=0$. Write $Z^n=\ker d$ (\emph{cocycles}), $B^n=\operatorname{im}d$ (\emph{coboundaries}), and $H^n(G,A)=Z^n/B^n$, with $B^0=0$.
\end{definition}

Unwinding: $H^0(G,A)=A^G$ is the group of fixed points. A $1$-cocycle is a \emph{crossed homomorphism}, $f(gh)=f(g)+g\cdot f(h)$, and the coboundaries are the \emph{principal} ones, $f(g)=g\cdot a-a$. If $G$ acts trivially, the
coboundaries are $0$ and $H^1(G,A)=\operatorname{Hom}(G,A)$. A $2$-cocycle satisfies
\[
g\cdot f(h,k)+f(g,hk)=f(gh,k)+f(g,h),
\]
an identity that will turn out to be associativity in disguise. A $2$-cocycle is \emph{normalised} if $f(1,g)=f(g,1)=0$ for all $g$.

\begin{lemma}\label{lem:normalised}
Every $2$-cocycle is cohomologous to a normalised one.
\end{lemma}

\begin{proof}
Put $g=h=1$ in the cocycle identity: $f(1,k)=f(1,1)$ for all $k$. Put $h=k=1$: $g\cdot f(1,1)=f(g,1)$. Let $b\in C^1$ be the constant function $b(g)=f(1,1)$; then $(db)(g,h)=g\cdot f(1,1)$. Now
$(f-db)(1,k)=f(1,1)-f(1,1)=0$ and $(f-db)(g,1)=g\cdot f(1,1)-g\cdot f(1,1)=0$.
\end{proof}

The first cohomology group has a concrete meaning for semidirect products. If $G$ acts on $A$, the \emph{semidirect product} $A\rtimes G$ is the set $A\times G$ with $(a,g)(b,h)=(a+g\cdot b,gh)$. A
\emph{complement} of $A=\{(a,1)\}$ is a subgroup $K$ with $K\cap A=\{0\}$ and $AK=A\rtimes G$; the subgroup $\{(0,g)\}$ is one.

\begin{proposition}\label{prop:complements}
The complements of $A$ in $A\rtimes G$ are the subgroups $K_f=\{(f(g),g):g\in G\}$ for $f\in Z^1(G,A)$, and $K_f,K_{f'}$ are conjugate by an element of $A$ if and only if $f-f'\in B^1$. So $H^1(G,A)$ is in bijection with the
$A$-conjugacy classes of complements.
\end{proposition}

\begin{proof}
A complement $K$ meets each coset $A\times\{g\}$ in exactly one element (if two, their quotient is in $K\cap A$; at least one, since $AK$ is everything), so $K=\{(f(g),g)\}$ for a function $f$. It is a subgroup if and only if
$(f(g),g)(f(h),h)=(f(g)+g\cdot f(h),gh)$ equals $(f(gh),gh)$, that is, if and only if $f$ is a crossed homomorphism; conversely every such $K_f$ is a complement. Conjugating by $(a,1)$:
\[
(a,1)(f(g),g)(-a,1)=\bigl(a+f(g)-g\cdot a,\ g\bigr),
\]
so $(a,1)K_f(a,1)^{-1}=K_{f-da}$.
\end{proof}

\section{Hilbert's Theorem 90}

\begin{theorem}[Hilbert, Noether, Speiser]\label{thm:hilbert90}
Let $L/K$ be a finite Galois extension with group $G$. Then $H^1(G,L^\times)=1$ and $H^1(G,L)=0$.
\end{theorem}

\begin{proof}
\emph{Multiplicative.} A $1$-cocycle is a map $f\colon G\to L^\times$ with $f(\sigma\tau)=f(\sigma)\,\sigma(f(\tau))$. The elements of $G$ are distinct automorphisms of $L$, so by Dedekind's lemma (Tome~2, Lemma~10.15, whose proof
works over any field) the map $\sum_{\tau}f(\tau)\tau\colon L\to L$ is not zero: choose $c$ with $b=\sum_\tau f(\tau)\tau(c)\ne0$. Then
\[
\sigma(b)=\sum_\tau\sigma(f(\tau))\,\sigma\tau(c)=\sum_\tau f(\sigma)^{-1}f(\sigma\tau)\,\sigma\tau(c)=f(\sigma)^{-1}b,
\]
since $\sigma\tau$ runs over $G$ with $\tau$. So $f(\sigma)=b\,\sigma(b)^{-1}$, a coboundary (written multiplicatively).

\emph{Additive.} The trace $\operatorname{Tr}=\sum_\tau\tau$ is not zero, again by Dedekind's lemma, and it is $K$-linear with values in $K$; choose $c$ with $\operatorname{Tr}(c)=1$. For a cocycle $f(\sigma\tau)=f(\sigma)+\sigma(f(\tau))$, put
$b=\sum_\tau f(\tau)\tau(c)$. Then
\[
\sigma(b)=\sum_\tau\bigl(f(\sigma\tau)-f(\sigma)\bigr)\sigma\tau(c)=b-f(\sigma)\operatorname{Tr}(c)=b-f(\sigma),
\]
so $f(\sigma)=b-\sigma(b)$.
\end{proof}

\begin{corollary}[Theorem~90 of the \emph{Zahlbericht}]\label{cor:hilbert90}
If $L/K$ is cyclic with generator $\sigma$ and $N(x)=x\,\sigma(x)\cdots\sigma^{n-1}(x)=1$, then $x=b/\sigma(b)$ for some $b\in L^\times$. If $\operatorname{Tr}(x)=0$, then $x=b-\sigma(b)$ for some $b\in L$.
\end{corollary}

\begin{proof}
Define $f(\sigma^i)=x\,\sigma(x)\cdots\sigma^{i-1}(x)$ for $0\le i<n$. Then $f(\sigma^i\sigma^j)=f(\sigma^i)\sigma^i(f(\sigma^j))$: for $i+j<n$ both sides are the product of $\sigma^k(x)$ for $k<i+j$, and for $i+j\ge n$ the right side is
$\prod_{k<i+j}\sigma^k(x)=N(x)\,\sigma^n\bigl(f(\sigma^{i+j-n})\bigr)=f(\sigma^{i+j-n})$, as $\sigma^n=\mathrm{id}$ and $N(x)=1$. So $f$ is a cocycle, and $x=f(\sigma)=b/\sigma(b)$ by the theorem. The additive case is the same with sums.
\end{proof}

\begin{figure}[ht]
\centering
\fig{ch07-circle}
\caption{Rational points on the unit circle, from Hilbert's Theorem~90 for $\Q(i)/\Q$: each point $\frac{(m+ni)^2}{m^2+n^2}$ with $1\le m,n\le7$ coprime is drawn, with the line from $(-1,0)$ of slope $\frac nm$ that finds it.}
\label{fig:rationalcircle}
\end{figure}

\begin{example}\label{ex:pythagoras}
For $\Q(i)/\Q$, $\sigma$ is complex conjugation and $N(a+bi)=a^2+b^2$. A rational point $(a,b)$ on the unit circle is an element of norm $1$, so $a+bi=\beta/\bar\beta$ with $\beta\in\Q(i)^\times$, and after clearing denominators
$\beta=m+ni$ with $m,n\in\Z$. Then
\[
a+bi=\frac{m+ni}{m-ni}=\frac{(m+ni)^2}{m^2+n^2}=\frac{m^2-n^2}{m^2+n^2}+\frac{2mn}{m^2+n^2}\,i.
\]
This is the classical parametrisation of Pythagorean triples $(m^2-n^2,2mn,m^2+n^2)$ (Tome~1, Chapter~9), now seen as the vanishing of a cohomology group (Figure~\ref{fig:rationalcircle}).
\end{example}

\section{Extensions and the second cohomology group}\label{sec:extensions}

Let $1\to A\xrightarrow{}E\xrightarrow{\pi}G\to1$ be an extension of groups with $A$ abelian (written additively inside $E$ when convenient). Then $G$ acts on $A$ by $g\cdot a=s\,a\,s^{-1}$ for any $s\in E$ with $\pi(s)=g$: another
choice $s'=bs$ with $b\in A$ gives the same result, as $A$ is abelian. Two extensions $E$, $E'$ of $G$ by the $G$-module $A$ are \emph{equivalent} if an isomorphism $E\to E'$ is the identity on $A$ and compatible with the maps to $G$.

\begin{theorem}[Schreier]\label{thm:h2ext}
Let $A$ be a $G$-module. The equivalence classes of extensions of $G$ by $A$ inducing the given action are in bijection with $H^2(G,A)$. The class $0$ corresponds to the semidirect product $A\rtimes G$.
\end{theorem}

\begin{proof}
\emph{From an extension to a class.} Choose a \emph{section} $s\colon G\to E$, $\pi s=\mathrm{id}$, with $s(1)=1$. Since $\pi(s(g)s(h))=\pi(s(gh))$, there is $f(g,h)\in A$ with
\[
s(g)s(h)=f(g,h)\,s(gh),
\]
a product in $E$. Expand $s(g)s(h)s(k)$ in two ways:
\[
\begin{gathered}
\bigl(s(g)s(h)\bigr)s(k)=f(g,h)s(gh)s(k)=f(g,h)f(gh,k)s(ghk),\\
s(g)\bigl(s(h)s(k)\bigr)=s(g)f(h,k)s(hk)=\bigl(g\cdot f(h,k)\bigr)s(g)s(hk)=\bigl(g\cdot f(h,k)\bigr)f(g,hk)s(ghk).
\end{gathered}
\]
In additive notation for $A$, $f(g,h)+f(gh,k)=g\cdot f(h,k)+f(g,hk)$: $f$ is a $2$-cocycle, normalised since $s(1)=1$. Another section is $s'(g)=b(g)s(g)$ with $b\colon G\to A$, $b(1)=0$, and then
\[
s'(g)s'(h)=b(g)\bigl(g\cdot b(h)\bigr)s(g)s(h)=b(g)\bigl(g\cdot b(h)\bigr)f(g,h)\,b(gh)^{-1}s'(gh),
\]
so $f'=f+db$. Equivalent extensions give the same cocycles for corresponding sections. So the class of $f$ in $H^2(G,A)$ depends only on the class of the extension.

\emph{From a cocycle to an extension.} By Lemma~\ref{lem:normalised}, every class contains a normalised cocycle $f$. Let $E_f=A\times G$ with
\[
(a,g)(b,h)=\bigl(a+g\cdot b+f(g,h),\ gh\bigr).
\]
Associativity: the first coordinates of the products $\bigl((a,g)(b,h)\bigr)(c,k)$ and
$(a,g)\bigl((b,h)(c,k)\bigr)$ are
\[
\begin{gathered}
a+g\cdot b+f(g,h)+gh\cdot c+f(gh,k),\\
a+g\cdot b+gh\cdot c+g\cdot f(h,k)+f(g,hk),
\end{gathered}
\]
equal by the cocycle identity. The identity element is
$(0,1)$ (by normalisation), and $\bigl(-g^{-1}\cdot a-f(g^{-1},g),\,g^{-1}\bigr)$ is a left inverse of $(a,g)$, which suffices. The maps $a\mapsto(a,1)$ and $(a,g)\mapsto g$ make $E_f$ an extension, inducing the given action since
$(0,g)(a,1)(0,g)^{-1}=(g\cdot a,1)$; with the section $s(g)=(0,g)$ its cocycle is $f$, because $(0,g)(0,h)=(f(g,h),gh)=(f(g,h),1)(0,gh)$.

\emph{Inverse bijections.} If $E$ has a section $s$ with cocycle $f$, the map $E_f\to E$, $(a,g)\mapsto a\,s(g)$, is bijective, compatible with the maps from $A$ and to $G$, and a homomorphism:
$a\,s(g)\,b\,s(h)=a\,(g\cdot b)\,s(g)s(h)=\bigl(a+g\cdot b+f(g,h)\bigr)s(gh)$. So every extension is equivalent to $E_f$ for each of its cocycles $f$. In particular, if $f'=f+db$ are normalised (then $b(1)=(db)(1,1)=0$), the extension $E_f$ has the
section $g\mapsto(b(g),g)$, whose cocycle is $f'$ by the first part, so $E_f$ is equivalent to $E_{f'}$. Thus extension classes and cohomology classes correspond bijectively.

Finally $f=0$ gives the multiplication $(a,g)(b,h)=(a+g\cdot b,gh)$ of $A\rtimes G$; and an extension has class $0$ if and only if it has a section that is a homomorphism, that is, if and only if it \emph{splits}.
\end{proof}

When $G$ acts trivially and $A$ is central, $H^2(G,A)$ classifies \emph{central} extensions. When moreover $G$ is abelian and only abelian $E$ are wanted, one restricts to symmetric cocycles, and gets back
$\operatorname{Ext}^1_\Z(G,A)$ of Chapter~6 (this can also be checked by enumeration).

\section{Cyclic groups, and carrying}

For a cyclic group everything is explicit. Let $G=\langle\sigma\rangle$ have order $n$, and let $N=1+\sigma+\dots+\sigma^{n-1}$ act on $A$; then $N(A)\subseteq A^G$, since $\sigma N=N$.

\begin{theorem}\label{thm:h2cyclic}
For $G$ cyclic of order $n$ and any $G$-module $A$, $H^2(G,A)\cong A^G/N(A)$. The class of $a\in A^G$ is represented by the cocycle
\[
f_a(\sigma^i,\sigma^j)=\begin{cases}a&\text{if }i+j\ge n,\\0&\text{if }i+j<n,\end{cases}\qquad0\le i,j<n.
\]
\end{theorem}

\begin{proof}
\emph{$f_a$ is a cocycle.} Let $g=\sigma^i$, $h=\sigma^j$, $k=\sigma^l$ with $0\le i,j,l<n$. Since $a$ is fixed by $G$, the identity $g\cdot f(h,k)+f(g,hk)=f(gh,k)+f(g,h)$ says that the number of ``carries'' is the same when the digits
$i,j,l$ are added modulo $n$ as $i+(j+l)$ or as $(i+j)+l$; both numbers equal $\lfloor(i+j+l)/n\rfloor$. It is normalised. In $E_{f_a}$, the element $t=(0,\sigma)$ satisfies $t^i=(0,\sigma^i)$ for $i<n$ (no carry yet) and
$t^n=t^{n-1}t=(f_a(\sigma^{n-1},\sigma),1)=(a,1)$.

\emph{An invariant of extensions.} Given an extension $1\to A\to E\to G\to1$, choose $t\in E$ mapping to $\sigma$. Then $t^n\in A$, and it is fixed by $G$, since $\sigma$ acts by conjugation by $t$, which commutes with $t^n$. Another choice is
$t'=bt$ with $b\in A$, and
\[
(bt)^n=b\cdot(tbt^{-1})\cdot(t^2bt^{-2})\cdots(t^{n-1}bt^{-(n-1)})\cdot t^n,
\]
that is, additively, $t'^n=N(b)+t^n$. So $\Phi(E)=[t^n]\in A^G/N(A)$ depends only on $E$ (and only on its equivalence class).

\emph{$\Phi$ is a bijection.} If $t^n=a$, the section $s(\sigma^i)=t^i$ ($0\le i<n$) has cocycle exactly $f_a$: $t^it^j$ is $t^{i+j}$ if $i+j<n$ and $a\,t^{i+j-n}$ otherwise. So $E$ is equivalent to $E_{f_a}$. Hence
$\Phi(E_{f_a})=[a]$ (so $\Phi$ is onto), and if $\Phi(E)=\Phi(E')$, choose $t$ in $E$ and $t'$ in $E'$ with $t^n=a$, $t'^n=a+N(b)$; replacing $t$ by $bt$ gives $E\simeq E_{f_{a+N(b)}}\simeq E'$, so $\Phi$ is injective.
Since $f_a+f_{a'}=f_{a+a'}$, the inverse map $[a]\mapsto[f_a]$ is a homomorphism.
\end{proof}

\begin{example}[carrying]\label{ex:carry}
Let $A=\Z/b$ with trivial action of $G=\Z/b$ (written additively, with generator $1$ and $\sigma^i=i$). Then $A^G=\Z/b$ and $N(x)=bx=0$, so $H^2(\Z/b,\Z/b)\cong\Z/b$. The cocycle $f_1(x,y)=1$ if $x+y\ge b$, $0$ otherwise, for digits
$0\le x,y<b$, is the carry of base-$b$ addition, and the extension it defines is $\Z/b^2$, with $(a,x)$ standing for the number $ab+x$. It generates $H^2$, and its multiples $f_a$, $a\in\Z/b$, give the $b$ classes. For $a=0$ the
extension is $\Z/b\times\Z/b$. In general $t=(0,1)$ has $t^b=(a,0)$, so $t$ has order $b$ times the order of $a$, and $E_{f_a}\cong\Z/b^2$ exactly when $a$ is a unit modulo $b$: there are $\varphi(b)$ inequivalent extensions with
middle group $\Z/b^2$, as in Example~\ref{ex:extzp}. For $b=10$ the other classes give $\Z/20\times\Z/5$ and $\Z/50\times\Z/2$ (Exercise~\ref{ex:carry10}).
\end{example}

\section{The Schur--Zassenhaus theorem}

\begin{proposition}\label{prop:killed}
Let $G$ be finite of order $m$. Then $m\cdot H^n(G,A)=0$ for $n=1,2$. If moreover $A$ is finite of order prime to $m$, then $H^1(G,A)=H^2(G,A)=0$.
\end{proposition}

\begin{proof}
For a $1$-cocycle, $f(hg)=f(h)+h\cdot f(g)$; summing over $g\in G$ and putting $b=\sum_gf(g)$ gives $b=mf(h)+h\cdot b$, so $mf(h)=b-h\cdot b=(d(-b))(h)$, a coboundary. For a $2$-cocycle, sum the identity
$g\cdot f(h,k)-f(gh,k)+f(g,hk)-f(g,h)=0$ over $k\in G$, and put $c(g)=\sum_kf(g,k)$: since $hk$ runs over $G$ with $k$, this gives $g\cdot c(h)-c(gh)+c(g)-mf(g,h)=0$, so $mf=dc$. Finally, if $|A|=r$ is prime to $m$, then
$H^n$ is killed by $r$ (since $A$ is) and by $m$, hence by $\gcd(r,m)=1$.
\end{proof}

\begin{theorem}[Schur, 1904; Zassenhaus, 1937]\label{thm:sz}
Let $E$ be a finite group and $A$ an abelian normal subgroup with $|A|$ prime to $|E/A|$. Then $A$ has a complement in $E$, and any two complements are conjugate.
\end{theorem}

\begin{proof}
By Proposition~\ref{prop:killed}, $H^2(E/A,A)=0$, so the extension splits (Theorem~\ref{thm:h2ext}): $E\cong A\rtimes G$ with $G=E/A$, and $A$ has a complement. Two complements $K$, $K'$ of $A$ are complements in $A\rtimes G$,
and $H^1(G,A)=0$ together with Proposition~\ref{prop:complements} shows that they are conjugate by an element of $A$.
\end{proof}

The theorem holds without assuming $A$ abelian; the existence of a complement follows by induction on $|E|$, using Sylow subgroups and the abelian case. The conjugacy of complements in general requires knowing that $A$ or
$E/A$ is solvable, and one of them has odd order since their orders are coprime; so the general conjugacy statement rests on the Feit--Thompson theorem of 1963, that groups of odd order are solvable, a 255-page proof.

\begin{beauty}
Hilbert's Theorem~90 found the rational points on a circle; in the hands of later mathematicians, $H^1(\operatorname{Gal}(L/K),\cdot)$ became the general tool for \emph{twisted forms}. Two objects defined over $K$ that become
isomorphic over $L$, like two quadratic forms or two algebras, differ by a $1$-cocycle with values in the automorphism group, and the cohomology class is exactly the obstruction to their being isomorphic over $K$. Theorem~90, in the form $H^1(G,\mathrm{GL}_n(L))=1$ that Speiser proved in 1919, says
that for vector spaces there is no obstruction at all: a vector space over $L$ with a compatible action of the Galois group comes from a vector space over $K$. This principle, \emph{Galois descent}, is how modern algebraic
geometry moves objects between fields, and $H^2(G,L^\times)$, the \emph{Brauer group}, classifies the central simple algebras over $K$ split by $L$, such as Hamilton's quaternions over $\R$, and is one of the pillars of class field theory.
\end{beauty}

\section{Exercises}

\begin{exercise}[$\star$]\label{ex:h1sign}
Let $G=\Z/2=\{1,\sigma\}$ act on $\Z$ by $\sigma\cdot a=-a$. Compute $H^0$, $H^1$ and $H^2$ (the latter with Theorem~\ref{thm:h2cyclic}), and do the same for the trivial action.
\end{exercise}

\begin{exercise}[$\star$]\label{ex:carry10}
In the setting of Example~\ref{ex:carry} with $b=10$, show that the extension with ``carry $a$'' is isomorphic, as a group, to $\Z/100$ when $\gcd(a,10)=1$, to $\Z/20\times\Z/5$ when $a=5$, to $\Z/50\times\Z/2$ when $a\in\{2,4,6,8\}$, and to
$(\Z/10)^2$ when $a=0$. How many inequivalent extensions have middle group $\Z/100$?
\end{exercise}

\begin{exercise}[$\star\star$]\label{ex:quaternion}
The quaternion group $Q_8=\{\pm1,\pm i,\pm j,\pm k\}$ is a central extension of $G=\{1,\bar i,\bar j,\bar k\}\cong(\Z/2)^2$ by $A=\{\pm1\}\cong\Z/2$. Compute its cocycle for the section $\bar g\mapsto g$ ($g\in\{1,i,j,k\}$), and show that it
is not a coboundary in two ways: directly, and by finding the complements.
\end{exercise}

\begin{exercise}[$\star\star$]\label{ex:cyclich1}
For $G=\langle\sigma\rangle$ cyclic of order $n$, show that $H^1(G,A)\cong\ker N/(\sigma-1)A$. Deduce Corollary~\ref{cor:hilbert90} again from Theorem~\ref{thm:hilbert90}, and conversely.
\end{exercise}

\begin{exercise}[$\star\star$]\label{ex:artinschreier}
(Artin--Schreier.) Let $K$ have characteristic $p$ and $L/K$ be cyclic of degree $p$ with generator $\sigma$. Show that there is $\alpha\in L$ with $\sigma(\alpha)=\alpha+1$, and that $\alpha$ is a root of $x^p-x-a$ with $a=\alpha^p-\alpha\in K$.
\end{exercise}

\begin{exercise}[$\star\star\star$]\label{ex:order20}
Use Theorem~\ref{thm:sz} to show that every group of order $20$ is a semidirect product $\Z/5\rtimes Q$ with $|Q|=4$, and that there are exactly five groups of order $20$ up to isomorphism.
\end{exercise}

\section{Hints}

\begin{hint}{ex:h1sign}
For the sign action, a crossed homomorphism is determined by $f(\sigma)$, with no condition. $N=1+\sigma$ is $0$ on $\Z$.
\end{hint}

\begin{hint}{ex:carry10}
The element $t=(0,1)$ satisfies $t^{10}=(a,0)$; compute its order and the order of the group it generates.
\end{hint}

\begin{hint}{ex:quaternion}
$ij=k$, $ji=-k$, $i^2=-1$. A complement would be a subgroup of order $4$ not containing $-1$.
\end{hint}

\begin{hint}{ex:cyclich1}
A crossed homomorphism is determined by $x=f(\sigma)$, and $f(\sigma^n)=f(1)=0$ is the condition $N(x)=0$.
\end{hint}

\begin{hint}{ex:artinschreier}
$\operatorname{Tr}(1)=p=0$. Apply the additive Theorem~90 to $x=-1$.
\end{hint}

\begin{hint}{ex:order20}
The Sylow $5$-subgroup is normal; count the homomorphisms $Q\to\operatorname{Aut}(\Z/5)\cong\Z/4$ up to automorphisms of $Q$.
\end{hint}

\section{Solutions}

\begin{solution}{ex:h1sign}
\idea For a cyclic group of order $2$, a crossed homomorphism is one value, and the norm is $1+\sigma$.
\writeup Sign action: $H^0=\{a:-a=a\}=0$. A crossed homomorphism has $f(1)=0$ and $f(\sigma)=x$, and the condition $f(\sigma^2)=f(\sigma)+\sigma f(\sigma)=x-x=0$ always holds; the principal ones are $\sigma\cdot b-b=-2b$. So
$H^1=\Z/2$. By Theorem~\ref{thm:h2cyclic}, $H^2=A^G/N(A)=0/0=0$. Trivial action: $H^0=\Z$, $H^1=\operatorname{Hom}(\Z/2,\Z)=0$, and $H^2=\Z/2\Z$ since $N(a)=2a$: the non-trivial class is the extension $\Z\xrightarrow2\Z$, that is
$0\to\Z\to\Z\to\Z/2\to0$, while the sign action's $H^1=\Z/2$ reflects the two conjugacy classes of complements $\{(0,1),(0,\sigma)\}$ and $\{(0,1),(1,\sigma)\}$ in the infinite dihedral group $\Z\rtimes\Z/2$.
\end{solution}

\begin{solution}{ex:carry10}
\idea The group is abelian, generated by $t=(0,1)$ and $A$, with $10t=(a,0)$; its isomorphism type follows from the order of $t$.
\writeup $E_a$ is abelian (the cocycle is symmetric and the action trivial), of order $100$, generated by $u=(1,0)$ and $t=(0,1)$, with relations $10u=0$, $10t=au$. If $\gcd(a,10)=1$, $t$ has order $10\cdot10=100$ and $E_a=\Z/100$. If $a=5$,
$t$ has order $20$, and $E_a$ has exponent $20$; since $|E_a|=100=20\cdot5$, it is $\Z/20\times\Z/5$ (the only abelian group of order $100$ and exponent $20$). If $a\in\{2,4,6,8\}$, $t$ has order $50$ and $E_a\cong\Z/50\times\Z/2$. If $a=0$,
$E_a=(\Z/10)^2$. So the $\varphi(10)=4$ classes $a=1,3,7,9$ give $\Z/100$: four inequivalent extensions with the same middle group.
\end{solution}

\begin{solution}{ex:quaternion}
\idea Compute $s(g)s(h)s(gh)^{-1}$ in a table; non-triviality shows up on the diagonal.
\writeup With $s(\bar g)=g$, the cocycle is $f(\bar g,\bar h)=s(\bar g)s(\bar h)s(\bar g\bar h)^{-1}$. For instance
\[
f(\bar i,\bar i)=i\cdot i=-1,\qquad f(\bar i,\bar j)=ij\,k^{-1}=1,\qquad f(\bar j,\bar i)=ji\,k^{-1}=-1,
\]
and likewise $f(\bar j,\bar j)=f(\bar k,\bar k)=-1$, with $f=1$ when an argument is $1$. In multiplicative notation with trivial action, a coboundary is $(db)(g,h)=b(g)b(h)b(gh)^{-1}$ with $b\colon G\to\{\pm1\}$. Then $(db)(1,1)=b(1)$ and $(db)(\bar i,\bar i)=b(\bar i)^2b(1)^{-1}=b(1)$, as $b(\bar i)^2=1$. If $f=db$,
then $b(1)=f(1,1)=1$ and so $f(\bar i,\bar i)=1$; but $f(\bar i,\bar i)=-1$. Contradiction. Alternatively, a complement would be a subgroup of order $4$ meeting $\{\pm1\}$ trivially; but every element of $Q_8$
other than $\pm1$ has order $4$ and square $-1$, so every subgroup of order $4$ contains $-1$.
\end{solution}

\begin{solution}{ex:cyclich1}
\idea A crossed homomorphism on a cyclic group is determined by its value on the generator, subject to one condition.
\writeup If $f$ is a crossed homomorphism and $x=f(\sigma)$, then $f(\sigma^i)=x+\sigma x+\dots+\sigma^{i-1}x$ by induction, and $0=f(1)=f(\sigma^n)=N(x)$. Conversely, for $x\in\ker N$ this formula defines a crossed homomorphism, as in the proof of
Corollary~\ref{cor:hilbert90}. Principal ones have $x=\sigma a-a$. So $H^1\cong\ker N/(\sigma-1)A$. For $A=L^\times$ this says: $H^1=1$ if and only if every element of norm $1$ is $b/\sigma(b)$, which is the corollary (with
$\sigma(b)/b$ or $b/\sigma(b)$, the same up to replacing $b$ by $b^{-1}$).
\end{solution}

\begin{solution}{ex:artinschreier}
\idea An element moved by $\sigma$ exactly by $+1$ is a root of an Artin--Schreier polynomial.
\writeup $\operatorname{Tr}(-1)=-p=0$, so by Corollary~\ref{cor:hilbert90} (additive) $-1=b-\sigma(b)$, that is $\sigma(b)=b+1$; put $\alpha=b$. Then $\sigma^i(\alpha)=\alpha+i$, so $\alpha\notin K$ and its conjugates
$\alpha,\alpha+1,\dots,\alpha+p-1$ are distinct: $L=K(\alpha)$ since $[L:K]=p$ is prime. Finally $\sigma(\alpha^p-\alpha)=(\alpha+1)^p-(\alpha+1)=\alpha^p+1-\alpha-1=\alpha^p-\alpha$ (Frobenius is additive in characteristic $p$), so
$a=\alpha^p-\alpha$ is fixed by $G$ and lies in $K$, and $\alpha$ is a root of $x^p-x-a$.
\end{solution}

\begin{solution}{ex:order20}
\idea The Sylow $5$-subgroup is normal of order prime to its index; Schur--Zassenhaus splits the group, and the semidirect products are classified by actions.
\writeup The number of Sylow $5$-subgroups divides $4$ and is $\equiv1\pmod5$, so it is $1$: $A\cong\Z/5$ is normal. As $\gcd(5,4)=1$, Theorem~\ref{thm:sz} gives a complement $Q$ of order $4$, so $E\cong\Z/5\rtimes_\rho Q$ with
$\rho\colon Q\to\operatorname{Aut}(\Z/5)\cong(\Z/5)^\times\cong\Z/4$. Semidirect products for $\rho$ and $\rho\circ\alpha$, $\alpha\in\operatorname{Aut}Q$, are isomorphic. For $Q=\Z/4$: $\rho$ trivial ($\Z/20$), of image of order $2$
(the generator acts by $-1$: the dicyclic group of order $20$), or injective (the two injective maps differ by the automorphism $-1$ of $\Z/4$: the Frobenius group $\Z/5\rtimes\Z/4$, the affine maps $x\mapsto ax+b$ of $\mathbb F_5$). For
$Q=(\Z/2)^2$: $\rho$ trivial ($\Z/10\times\Z/2$) or with image $\{\pm1\}$ (all three such maps are permuted by $\operatorname{Aut}Q$: the dihedral group of order $20$). These five groups are pairwise non-isomorphic. The Sylow $2$-subgroup of $\Z/5\rtimes Q$ is isomorphic to $Q$. The two abelian ones differ by it, and so does the dihedral group from the other two non-abelian ones; and the dicyclic group has
centre of order $2$ (the square of the generator of $Q$ acts trivially), while the Frobenius group has trivial centre. So there are exactly five.
\end{solution}

\section*{Sources and further reading}
\addcontentsline{toc}{section}{Sources and further reading}

D.~C.~Isaksen, ``A cohomological viewpoint on elementary school arithmetic'', \emph{Amer.\ Math.\ Monthly} 109 (2002), 796--805; D.~Hilbert, ``Die Theorie der algebraischen Zahlkörper'', \emph{Jahresber.\ Deutsch.\ Math.-Verein.} 4
(1897), 175--546, Satz~90; O.~Hölder, ``Die Gruppen der Ordnungen $p^3$, $pq^2$, $pqr$, $p^4$'', \emph{Math.\ Ann.} 43 (1893), 301--412; O.~Schreier, ``Über die Erweiterung von Gruppen I'', \emph{Monatsh.\ Math.\ Phys.} 34 (1926),
165--180; I.~Schur, ``Über die Darstellung der endlichen Gruppen durch gebrochene lineare Substitutionen'', \emph{J.\ Reine Angew.\ Math.} 127 (1904), 20--50; H.~Zassenhaus, \emph{Lehrbuch der Gruppentheorie} (1937); W.~Feit and
J.~G.~Thompson, ``Solvability of groups of odd order'', \emph{Pacific J.\ Math.} 13 (1963), 775--1029.

For more: K.~S.~Brown, \emph{Cohomology of Groups} (Springer, 1982), Chapters~III and~IV; J.-P.~Serre, \emph{Local Fields} (Springer, 1979), Chapters~VII and~X, for Galois cohomology and Theorem~90; C.~A.~Weibel,
\emph{An Introduction to Homological Algebra}, Chapter~6.
